\honest{Positive Bias: Look Into the Dark}{Eliezer Yudkowsky}{2007}

I write 2-4-6 on the board and say that it fits a rule; you may test other triplets to find the rule. One student tests 4-6-2 (No), 4-6-8 (Yes) and 10-12-14 (Yes), then writes down a guess. The real rule is any three numbers in ascending order.

The task is Peter Wason's. Subjects are usually confident, yet only about a fifth find the rule, in the original and in replications. Most test triplets that their own guess predicts will fit, hear ``Yes,'' and announce the guess. Someone who guesses ``increasing by two'' tests 8-10-12. Someone who guesses X, 2X, 3X tests 3-6-9. Some announce rules far more complicated than the real one. To find the real rule, you must test triplets that should not fit, such as 20-23-26.

My own student's first test was 4-6-2, a triplet that ``increasing by two'' forbids. It came back No. \nb{It could not tell the student's guess from the true rule, since both forbid it. The hard part is choosing which case that should not fit to test, and the post does not say how to choose.}

\nb{Wason designed the task so that, in his words, confirming evidence alone ``will almost certainly lead to erroneous conclusions'': the true rule is much broader than any natural guess. Twenty years before the post, Klayman and Ha showed that testing the cases you expect to fit is often the better strategy, and that it fails mainly in tasks built like this one. The post does not mention them.}

It seems to me that testing positive cases should be distinguished from defending the belief you started with. \nb{Klayman and Ha drew that distinction in 1987.}

Phlogiston is my next example. It could explain a flame going out, and could just as well have explained it not going out.

Looking for negative cases goes against what experiment has shown to be ``human instinct.'' ``For by instinct, we human beings only live in half the world.''

The 2-4-6 task is cold logic, not emotional, and still people fail. The mistake is below the level of words, in which example springs to mind. So knowing in words that you should test both kinds will not by itself fix it. You have to learn, wordlessly, ``to flinch toward the zero, instead of away from it.'' \nb{Later research, reported by Gale and Ball in 2006, found that rewording the task, asking subjects to find two complementary rules instead of one, raised success from about a fifth to over 60 per cent. That does not show that teaching the rule works, but it shows the error depends on how the task is framed, not only on instinct. The post reports only the fifth.} I remind you that I have long argued that ``the strength of a hypothesis is what it can't explain, not what it can.''

So to judge an explanation, look for results it could not explain. A commenter offered superconductivity and ferromagnetism as examples of emergence. I replied that their absence would be examples of emergence too, which was the problem. I confess that I too fell for the 2-4-6 task the first time I read about it, and that ``I'm still working on it myself.''

So much of this skill is below words that people agree with you and then, in the next sentence, do the opposite. I write partly to see what my words have not conveyed. Are you, right now, searching only for positive examples of positive bias?
